% =====================================================================
% Section 5. Conclusions
% =====================================================================
\section{Conclusions}
\label{sec:conclusions}

We have presented a second-order semi-implicit IMEX finite volume scheme for the single-velocity six-equation
two-phase model with total energy and mechanical, thermal and chemical relaxation, for stiffened-gas,
Noble--Abel stiffened-gas and Mie--Gr\"uneisen equations of state, in one and two space dimensions.

The two stiffnesses of the model, the acoustic waves at low Mach numbers and the relaxation processes up to infinite
rates, are treated in a single implicit step. The implicit pressure work and the three relaxation processes, each at a
finite or infinite rate, a slower process being active only when the faster ones are instantaneous \cite{Pelanti2022}, enter one local closure per cell, which reduces to a scalar equation for the
cell pressure; each stage then reduces to a mildly nonlinear system for the cell pressures with a three-point
(five-point in 2D) Jacobian, solved by Newton's method. The time step is restricted by the flow velocity only, whatever
the Mach number and the relaxation rates. We have proved that, under non-degeneracy assumptions, the scheme converges
to the corresponding equilibrium scheme as each relaxation rate tends to infinity and reduces to the scheme without that
relaxation as the rate vanishes; formally, it is asymptotic preserving in the low Mach number limit of the mixture.

At shocks, the difference between the conservative and the non-conservative mixture internal energy is distributed
between the phases according to the shock relations of the model: the physical part with the constant share for which
the shock profile of the relaxed model ends on the Hugoniot of each phase, computed by shooting, which for weak shocks
reduces to the unique consistent local share, and the discretisation error with the unique share that keeps both phases
admissible for every defect. No calibrated heat exchange is needed.

The flux is split following Zha and Bilgen: all quantities are transported explicitly with the old velocity by the
path-conservative Rusanov scheme, whose dissipation speed is the material speed, the volume fraction being advected
entirely in this step, while the pressure, its work on each phase and an incremental flux of the conservative transported
quantities are treated implicitly. The central implicit operator is supplemented by implicit Laplace dissipation of the
pressure and the velocity with acoustic coefficients, weighted by the deficit of the implicit viscosity with respect to
the upwind one, so that it vanishes in the low Mach number regime, and limited to second order in smooth regions. Second
order in time is obtained with the stiffly accurate, L-stable LSDIRK2 scheme. We have also proved that the scheme
preserves moving contacts and the phase temperatures under transport.

The numerical results confirm these properties: second-order convergence on a smooth two-phase vortex, a Newton
iteration count and a time step independent of the Mach number, face velocities that are discretely divergence-free to
$O(M^2)$ down to $M\approx10^{-6}$ on the Taylor--Green vortex, and agreement with exact solutions and reference results
for interfaces, single-phase and two-phase shock tubes, cavitation, evaporation and depressurisation problems, and
two-dimensional shock--bubble, Richtmyer--Meshkov, bubble-collapse and underwater-explosion problems, at instantaneous
and finite relaxation rates.

Future work will address an extension to higher order in space and time, unstructured meshes, and a velocity
disequilibrium model, in which the drag relaxation would enter the same local closure.
